% Job posting: https://clarius.com/career/co/research-development-rd/0D.17C/machine-learning-engineer/all/
% =====================================================================
%  CV - Nishal Stanislaus Silva, PhD
%  Target: Clarius Mobile Health - Machine Learning Engineer (R&D, 24-month project)
%  Location: Remote, Canada (Vancouver BC hybrid preferred)
%  Variant: V5 (Computer Vision / Industrial)
% =====================================================================

\documentclass[11pt]{article}
\usepackage[left=0.7in,top=0.7in,right=0.7in,bottom=0.7in]{geometry}
\usepackage{enumitem}
\usepackage[hidelinks]{hyperref}
\usepackage{titlesec}
\usepackage{array}
\usepackage{setspace}
\usepackage{needspace}
\usepackage[T1]{fontenc}
\usepackage{newtxtext,newtxmath}
\ifdefined\pdfgentounicode
  \input{glyphtounicode}
  \pdfgentounicode=1
\fi
\setstretch{1.05}
\pagenumbering{gobble}
\titleformat{\section}{\Large\scshape}{}{4em}{}[\vspace{-0.7em}\rule{\linewidth}{0.4pt}\vspace{-0.4em}]
\titlespacing*{\section}{0pt}{1em}{0.4em}
\newenvironment{cvrole}[5]{%
  \par\noindent\textbf{\large #1} | \textit{\small #2, #3}\hfill \textit{\small #4 - #5}\par
  \begin{itemize}[leftmargin=2em, itemsep=-0.3em, topsep=-0.5em]%
}{%
  \end{itemize}\par\noindent\mbox{}\par\vspace{0.1em}%
}
\newcommand{\cvName}{Nishal Stanislaus Silva}
\newcommand{\cvEmail}{nishal.silva@hotmail.com}
\newcommand{\cvPhone}{+1 613 200 2030}
\newcommand{\cvCity}{Toronto, ON}
\newcommand{\cvSite}{https://nishal.xyz}
\newcommand{\cvLinkedIn}{https://www.linkedin.com/in/nishal-silva}
\newcommand{\cvGitHub}{https://github.com/ns2max}
\newcommand{\cvStatus}{Canadian Permanent Resident, Eligible to work in Canada without sponsorship}
\newcommand{\cvTagline}{Computer Vision \& ML Engineer | Industrial Inspection, IoT \& Edge Deployment}

\begin{document}
\pagestyle{empty}
\begin{center}
    {\LARGE \scshape \textbf{\cvName}}{\large \textit{, PhD}}\\[1pt]
    {\small \cvTagline}\\[1pt]
    {\footnotesize\textit{\cvStatus}}\\[1pt]
    {\small
    \href{mailto:\cvEmail}{\cvEmail} \,\,|\,\, \cvPhone \,\,|\,\, \cvCity
    \,\,|\,\, \href{\cvSite}{nishal.xyz} \,\,|\,\, \href{\cvLinkedIn}{linkedin.com/in/nishal-silva}
    \,\,|\,\, \href{\cvGitHub}{github.com/ns2max}}
\end{center}

\section*{Professional Summary}

Computer vision and ML engineer with a PhD (University of Trento, 2025) and 10+ years across industry and academia, taking ML models from research through production on real hardware. At MAS Holdings I built and deployed production computer-vision inspection for apparel manufacturing at scale, feeding camera, sensor and IoT streams into live workflows and cutting inspection time 99.5\%, with CI/CD, code review and pytest that I introduced. My doctoral work carried the same engineering discipline to the model side: a real-time detector running inference in 14 ms on a Raspberry Pi 4 at 31.4\% CPU, F1 0.76, 74$\times$ faster than the DTW baseline it replaced (JAES 2026). Optimizing model performance and inference speed for deployment is my core strength. My closest health work is a COVID-19 remote-vitals computer-vision device deployed across hospital wards (GMOA recognition); ultrasound imaging is a domain I would ramp on, with a record of learning applied fields cold.

\section*{Core Competencies}

\textbf{Deep Learning \& Computer Vision:} deep learning for vision and sensor data, production OpenCV inspection pipelines, template and feature matching, image registration and geometric warping, OCR and layout segmentation, explainable defect overlays, industrial and microscopic camera arrays, CNN and sequence architectures (RNN, LSTM, GRU), similarity-based and open-set matching\\
\textbf{ML Engineering:} research-to-production model delivery, training and evaluation under domain shift, automated model retraining, evaluation and monitoring, failure diagnosis and root-cause analysis, benchmark and dataset design, statistical validation, metrics tied to operational outcomes, TensorFlow/Keras, PyTorch, scikit-learn\\
\textbf{Systems \& Deployment:} inference-speed and model-performance optimization for production, real-time inference under sub-30 ms budgets, ARM and Raspberry Pi latency, memory and CPU profiling, embedded Linux, C++17 real-time engines, Python on Unix/Linux, ML data pipelines, Docker, AWS, version control, testing and code review (Git, pytest, CI/CD), ML technical-debt reduction

\section*{Technical Stack}

\textbf{Languages:} Python (with Unix/Linux), C++17, C, JavaScript/TypeScript, SQL, MATLAB, Shell\\
\textbf{ML \& AI:} TensorFlow, Keras, PyTorch, scikit-learn, NumPy, Pandas, SciPy; RNN/CNN, embedding and similarity matching, VAE and diffusion prototypes; model compression, quantization, ONNX and TFLite familiarity; ARM latency profiling\\
\textbf{Computer Vision:} OpenCV, template and feature matching, OCR, image registration and warping, corner detection, explainable overlays, industrial and microscopic cameras, projector-camera AR\\
\textbf{Data \& Dev:} AWS (EC2, S3, ECS, MWAA, RDS), ETL and real-time ingestion pipelines, Docker, Airflow, Jenkins, Git, Linux/UNIX, pytest, CI/CD, REST APIs; Kubernetes and GCP familiarity

\section*{Education}
\textbf{Ph.D - Information Engineering and Computer Science} \textit{\small{ - University of Trento, Italy}} \hfill \textit{Jan 2025}\\[2pt]
\noindent\textit{\small Thesis: Embedded Real-time Musical Pattern Detection for Smart Musical Instruments}\\[3pt]
\noindent\textbf{M.Sc - Telecommunication and Electronic Engineering} \textit{\small{ - Sheffield Hallam University, UK}} \hfill \textit{Aug 2018}\\[3pt]
\noindent\textbf{B.Eng (Hons) - Electronic Engineering} \textit{\small{ - Sheffield Hallam University, UK}} \hfill \textit{Mar 2014}

\section*{Professional Experience}

\begin{cvrole}{Research Engineer}{MAS Holdings (Pvt) Ltd}{Colombo, Sri Lanka}{Jan 2015}{Jul 2020}
    \item Built, trained and deployed end-to-end computer-vision and sensor systems for apparel manufacturing at scale, taking model-based inspection from prototype to production on the factory floor; operational efficiency improved by up to 300\%.
    \item Shipped automated multilingual care-label QC: OpenCV template and feature matching with an explainable defect overlay, iterated from scanner to industrial camera to a conveyor with PLC reject; inspection time cut 99.5\%.
    \item Built loom-side fabric defect detection with a microscopic camera array (overlapping fields of view across 60-inch rolls, 5 to 6 second alert budget); pilot cut operator headcount 50\%.
    \item Delivered the Promptly on-demand garment-printing system: designed the camera-to-warp geometry-alignment algorithm, personally built the RIP integration into the automation loop, and consulted on the flipping mechanism; now a Twinery/MAS commercial product across four countries.
    \item Built real-time ingestion and ETL for high-frequency IoT and operational time-series across sites in South Asia, North America and Africa; ran C++/Python inference on embedded hardware; introduced CI/CD, code review and pytest; mentored engineers and delivered group-wide CV training.
\end{cvrole}

\needspace{5\baselineskip}
\begin{cvrole}{Postdoctoral Researcher}{University of Trento}{Trento, Italy}{Jan 2025}{Dec 2025}
    \item Owned the full ML pipeline for streaming audio and IMU/gesture data: acquisition, feature engineering, data labeling, training and evaluation under domain shift, edge deployment on Raspberry Pi 4 with production monitoring; EU Horizon EIC Pathfinder project MUSMET.
    \item Delivered real-time inference at 14 ms on a Raspberry Pi 4, F1 0.76, 31.4\% CPU, 74$\times$ faster than the 1038 ms DTW baseline (JAES 2026); designed the frozen-protocol benchmark suites and ablations that quantified the accuracy, latency and CPU trade-offs.
\end{cvrole}

\needspace{5\baselineskip}
\begin{cvrole}{Visiting Researcher}{McGill University (IDMIL / CIRMMT)}{Montreal, Canada}{May 2024}{Aug 2024}
    \item Fused inertial (IMU) and audio streams through a hierarchical finite state machine, accumulating repeated observations of the same event to reach F1 0.78 across a 500-event corpus (IEEE IS\textsuperscript{2} 2025); profiled compute and power for real-time feasibility on self-powered embedded hardware.
\end{cvrole}

\section*{Selected Projects \footnotesize {\normalfont{(full portfolio: \href{https://nishal.xyz/research}{\textit{nishal.xyz/research}})}}}
\begin{itemize}[leftmargin=1.2em, itemsep=-0.25em]
    \item \textbf{COVID-19 remote vitals device (MAS):} computer-vision system digitizing patient vitals from retrofitted webcams into a server dashboard, deployed across multiple hospital wards for remote monitoring; recognized by the Government Medical Officers' Association of Sri Lanka.
    \item \textbf{Real-time audio pattern detection:} MFCC front-end into a stacked RNN on 30 ms windows; Raspberry Pi 4 at 14 ms, 31.4\% CPU, F1 0.76 against a DTW baseline at 1038 ms.
    \item \textbf{SSIM-based training-free detection:} four-attribute representation of symbolic music with bicubic resizing for variable length; 95\% detection across human, synthetic and JKUPDD sets with no training data.
    \item \textbf{Tech-pack digitization and lace archive (MAS, Noyon Dentelle):} pre-LLM OCR and layout segmentation of physical blueprints; flatbed and microscopic capture with context-aware ERP search, still in production use.
\end{itemize}

\enlargethispage{3\baselineskip}
\vspace{-0.5em}
\section*{Publications \& Recognition}
\begin{itemize}[leftmargin=1.2em, itemsep=-0.25em]
    \item Silva, N. \& Turchet, L. (2026). Real-Time Audio Pattern Detection for Smart Musical Instruments. \textit{J. Audio Eng. Soc.} 74(3).
    \item 10+ peer-reviewed papers (JAES, IEEE IS\textsuperscript{2}, IEEE I3DA, ACM Audio Mostly, DAFx, IEEE Asilomar); four open-access Zenodo datasets (9,800+ recordings, 70+ musicians); MIDI Innovation Awards 2023 Finalist.
    \item Nebula, a C++17 audio-feature-extraction library with per-feature ARM/Raspberry Pi latency benchmarks; recent side projects built with AI-assisted workflows including Claude Code.
\end{itemize}

\end{document}
